%=============================================================================!
% 7.2  HAMILTON'S EQUATIONS                                                   !
%=============================================================================!
\section{Hamilton's equations}
\label{sec:ha-equations}

The Hamiltonian is built from the Lagrangian, so the Euler--Lagrange
equation must be expressible through it. Finding how needs the partial
derivatives of $H$, and these differ from those of $\Lag$ in what is held
fixed: $\partial\Lag/\partial q$ is taken at fixed $\dot{q}$, while
$\partial H/\partial q$ is taken at fixed $p$, and at fixed $p$ the rate
$\dot{q}$ changes with $q$. The chain rule of \cref{sec:en-gradient}
keeps track of this once it is stated for functions of several
variables that themselves depend on several others.

\begin{toolbox}[The chain rule in general]
Let $f(x, y)$ depend on two variables, and let these depend in turn on
two others, $x(s, t)$ and $y(s, t)$. Holding $s$ fixed and letting $t$
change moves the point $(x, y)$ along a path, so the chain rule along a
path (\cref{sec:en-gradient}) gives
\begin{equation*}
   \frac{\partial f}{\partial t} = \frac{\partial f}{\partial x}\,
   \frac{\partial x}{\partial t} + \frac{\partial f}{\partial y}\,
   \frac{\partial y}{\partial t} ,
\end{equation*}
where every derivative with respect to $t$ is taken at fixed $s$, and
$\partial f/\partial x$ and $\partial f/\partial y$ are evaluated at the
point $\bigl(x(s, t), y(s, t)\bigr)$. Holding $t$ fixed instead gives
the same with $s$ in place of $t$, and with more variables there is one
term for each variable of $f$. For example, $f = x^2y$ with $x = s + t$
and $y = s - t$ has
\begin{equation*}
   \frac{\partial f}{\partial t} = 2xy\times1 + x^2\times(-1)
   = 2(s + t)(s - t) - (s + t)^2 ,
\end{equation*}
which is what differentiating $f = (s + t)^2(s - t)$ directly with respect
to $t$ gives, by the product rule. A partial derivative is only defined
once we say which variables are held fixed: $\partial f/\partial x$ at
fixed $y$ is $2xy$, but if $y$ were written as $s - t$ and $x$ as $s + t$
and we asked how $f$ changes with $x$ at fixed $s$, so that $y = 2s - x$
changes too, then $f = x^2(2s - x)$ and the answer would be $2x(2s - x) -
x^2 = 2xy - x^2$, which differs from $2xy$.
\end{toolbox}

\subsection*{The equations}

In \cref{eq:ha-hamiltonian}, $\dot{q}$ is a function of $q$ and $p$,
found from $p = \partial\Lag/\partial\dot{q}$, so
\begin{equation*}
   H(q, p) = p\,\dot{q}(q, p) - \Lag\bigl(q, \dot{q}(q, p)\bigr) .
\end{equation*}
The variable $p$ enters through the factor $p$ and through $\dot{q}$, and
$\Lag$ depends on $p$ only through its second variable $\dot{q}$. At fixed
$q$, by the product rule for the first term and the chain rule for the
second,
\begin{align*}
   \frac{\partial H}{\partial p}
   &= \dot{q} + p\,\frac{\partial\dot{q}}{\partial p}
   - \frac{\partial\Lag}{\partial\dot{q}}\,\frac{\partial\dot{q}}
   {\partial p}
   && \text{differentiating at fixed $q$}\\
   &= \dot{q}
   && \text{since $\partial\Lag/\partial\dot{q} = p$.}
\end{align*}
The variable $q$ enters $p\,\dot{q}$ through $\dot{q}(q, p)$, and $\Lag$
twice, as its first variable, whose rate of change with $q$ is 1, and
through $\dot{q}(q, p)$. At fixed $p$, the general chain rule, with $\Lag$
for $f$, $q$ and $\dot{q}$ for $x$ and $y$, and $q$ for $t$, gives one
term for each,
\begin{align*}
   \frac{\partial H}{\partial q}
   &= p\,\frac{\partial\dot{q}}{\partial q}
   - \frac{\partial\Lag}{\partial q}
   - \frac{\partial\Lag}{\partial\dot{q}}\,\frac{\partial\dot{q}}
   {\partial q}
   && \text{differentiating at fixed $p$}\\
   &= -\frac{\partial\Lag}{\partial q}
   && \text{since $\partial\Lag/\partial\dot{q} = p$ again,}
\end{align*}
where $\partial\Lag/\partial q$ is the derivative of \cref{sec:la-euler},
at fixed $\dot{q}$. The Euler--Lagrange equation \cref{eq:la-euler} says
$\dot{p} = \dd(\partial\Lag/\partial\dot{q})/\dd t = \partial\Lag/
\partial q$, which is $-\partial H/\partial q$. Together,
\begin{equation}
   \dot{q} = \frac{\partial H}{\partial p} ,
   \qquad
   \dot{p} = -\frac{\partial H}{\partial q} .
   \label{eq:ha-hamilton}
\end{equation}
These are Hamilton's equations, one pair for each coordinate, found in
the same way for each $k$ when there are several. In place of one
equation in the second derivative $\ddot{q}$ there are two in first
derivatives, so the state $(q, p)$ at one instant fixes the motion, as
\cref{sec:ne-equations} found for $(x, v)$.

For the block, $H = p^2/(2m) + \frac12 kx^2$ and \cref{eq:ha-hamilton} give $\dot{x} = p/m$ and
$\dot{p} = -kx$, so $m\ddot{x} = \dot{p} = -kx$, the spring of
\cref{sec:ne-spring}. For the pendulum, $\dot\theta = p/(mL^2)$ and
$\dot{p} = -mgL\sin\theta$, which together give $mL^2\ddot\theta = \dot{p}
= -mgL\sin\theta$, \cref{eq:la-pendulum}.
For $N$ atoms, $\dot{\vb{r}}_i = \vb{p}_i/m_i$ and $\dot{\vb{p}}_i =
-\grad_iU = \vb{F}_i$: Newton's second law in the first-order form of
\cref{eq:ne-first-order}, with the momentum in place of the velocity.
\Cref{ex:ha-bead} recovers the bead's equation \cref{eq:la-bead} from its
Hamiltonian.

\subsection*{Phase space and the flow}

For one coordinate the states $(q, p)$ fill a plane, the phase plane of
\cref{sec:os-phase} with the momentum in place of the velocity; for $n$
coordinates they fill a space of $2n$ dimensions, and for $N$ atoms one
of $6N$. This is phase space. At each of its points Hamilton's equations
give the velocity of the state, $(\dot{q}, \dot{p}) = (\partial H/
\partial p, -\partial H/\partial q)$, a velocity attached to every point,
as the current of a river has a velocity at every point of its surface.
A state moves as a leaf carried by such a current, and the motions of
the system are the paths of the leaves, which, by the uniqueness of
\cref{sec:ne-equations}, never cross when $H$ does not contain the time;
the separatrix of the pendulum below only seems to cross itself at the
upright states, which are motions of their own (\cref{sec:os-phase}). This motion of all states at once is
called the Hamiltonian flow.

The flow keeps the energy. Along any motion, by the chain rule along a
path,
\begin{equation*}
   \frac{\dd H}{\dd t} = \frac{\partial H}{\partial q}\,\dot{q}
   + \frac{\partial H}{\partial p}\,\dot{p}
   = \frac{\partial H}{\partial q}\,\frac{\partial H}{\partial p}
   - \frac{\partial H}{\partial p}\,\frac{\partial H}{\partial q} = 0 .
\end{equation*}
Equivalently, the velocity of the flow is at right angles to the
gradient $(\partial H/\partial q, \partial H/\partial p)$, since their
dot product is the same zero, and so runs along the contours of $H$
(\cref{sec:en-gradient}). The phase portrait is therefore the contour
map of $H$. When $H$ contains the time itself, the chain rule has a third
term, and $\dd H/\dd t = \partial H/\partial t$. The block pushed by the
force $F_0\cos\omega t$ of \cref{sec:os-driven}, without drag, has $U =
\frac12 kx^2 - F_0x\cos\omega t$, whose slope gives the spring's force and
the push, and $\partial H/\partial t = \omega F_0x\sin\omega t$, which is not zero:
$H$, which now includes the push's term $-F_0x\cos\omega t$, rises and
falls as the push feeds energy in and takes it out.

\subsection*{The pendulum's portrait}

For the pendulum of \cref{sec:la-multipliers}, a bob of \SI{0.5}{\kilogram} on a
rod of \SI{1.2}{\metre}, $mL^2 = \SI{0.72}{\kilogram\metre\squared}$ and
$mgL = \SI{5.884}{\joule}$. The contours of $H$ (\cref{fig:ha-portrait}a)
are of three kinds. Below the energy $2mgL$ of the bob held upright, they
are closed loops around $\theta = 0$: the bob swings, and $p$ changes
sign at each end of the swing. Above $2mgL$ they are open, wavy lines on
which $p$ never vanishes: the bob goes over the top. The contour at
$2mgL$ itself passes through the upright states $(\pm\pi, 0)$ and
separates the two kinds; it is called the separatrix. On it, at the
bottom, $p^2/(2mL^2) = 2mgL$, so
\begin{equation*}
   p = \sqrt{4m^2gL^3} = 2mL^2\sqrt{\frac{g}{L}} = 2mL^2\omega_0
   = 2\times0.72\times2.8587 = \SI{4.117}{\kilogram\metre\squared\per
   \second} ,
\end{equation*}
with $\omega_0 = \sqrt{g/L}$ of \cref{sec:os-small}.

\begin{figure}[htb]
\centering
\includegraphics{ch07_hamilton/portrait}
\caption{The pendulum of \SI{0.5}{\kilogram} on a rod of
\SI{1.2}{\metre}. (a) Contours of $H$ in the phase plane, with arrowheads
showing the direction of the flow: swinging, below the upright energy
$2mgL$ (grey); the separatrix, at $2mgL$ (teal); going over the top
(ochre). (b) Started at the bottom with the momentum of the separatrix,
the angle still to go, $\pi - \theta$, from Hamilton's equations (dots),
on a logarithmic axis, against a line that falls as
$\mathrm{e}^{-\omega_0t}$ (dashed).}
\label{fig:ha-portrait}
\end{figure}

A bob given exactly this momentum at the bottom creeps towards the top
and never arrives, as \cref{sec:os-phase} said. On the separatrix the
energy is $2mgL$ at every angle, $\frac12 mL^2\dot\theta^2 + mgL(1 -
\cos\theta) = 2mgL$, so
\begin{align*}
   \dot\theta^2 &= \frac{2g}{L}\,(1 + \cos\theta)
   && \text{subtracting $mgL(1 - \cos\theta)$, dividing by
   $\tfrac12 mL^2$}\\
   &= \frac{4g}{L}\cos^2\frac{\theta}{2}
   && \text{by $\cos\theta = 2\cos^2\tfrac{\theta}{2} - 1$}\\
   \dot\theta &= 2\omega_0\cos\frac{\theta}{2}
   && \text{the root for motion towards $+\pi$,}
\end{align*}
taking the positive root because $\cos\frac{\theta}{2} > 0$ for $-\pi <
\theta < \pi$, and where the second line uses the double-angle formula of
\cref{sec:mo-trig},
$\cos\theta = \cos^2\frac{\theta}{2} - \sin^2\frac{\theta}{2}$, with
$\sin^2 = 1 - \cos^2$. Near the top, we write $\varphi = \pi - \theta$
for the angle still to go. Then $\cos\frac{\theta}{2} = \cos(\frac{\pi}{2} -
\frac{\varphi}{2}) = \sin\frac{\varphi}{2}$ (\cref{sec:mo-trig}), and $\sin\frac{\varphi}{2}
\approx \frac{\varphi}{2}$ for small $\varphi$ (\cref{sec:mo-taylor}), and $\dot\varphi = -\dot\theta$,
so
\begin{equation*}
   \dot\varphi \approx -\omega_0\,\varphi .
\end{equation*}
This is the equation of the drag of \cref{sec:ne-drag} with $\omega_0$ in
place of $\gamma$ and nothing to approach but zero, so $\varphi$ falls as
$\mathrm{e}^{-\omega_0t}$: it halves every $\ln 2/\omega_0 =
\SI{0.2425}{\second}$, however small it is, and never reaches zero.
Computed from Hamilton's equations (\cref{fig:ha-portrait}b), the bob
started at the bottom is \SI{0.229}{\radian} short of the top after
\SI{1}{\second} and \SI{0.0132}{\radian}, \ang{0.75}, short after
\SI{2}{\second}, and from then on the gap halves every
\SI{0.2425}{\second}.

\begin{keyidea}
Hamilton's equations, $\dot{q} = \partial H/\partial p$ and $\dot{p} =
-\partial H/\partial q$, move every state of phase space along a flow
that runs along the contours of $H$, so $H$ is conserved unless it
contains the time.
\end{keyidea}

\notebook{07}{launch a pendulum anywhere on its phase portrait}

\begin{selfcheck}
A ball thrown straight up has $H = p^2/(2m) + mgq$, with $q$ its height.
Write Hamilton's equations for it, and describe the contours of $H$ in
its phase plane.
\end{selfcheck}
